\documentclass[11pt]{article}
\usepackage[margin=25mm]{geometry}
\usepackage{fontspec}
\setmainfont{Latin Modern Roman}
\usepackage{amsmath,amssymb,mathtools}
\usepackage{unicode-math}
\setmathfont{Latin Modern Math}
\usepackage{xurl,hyperref}
\hypersetup{colorlinks=true,urlcolor=blue}
\setcounter{secnumdepth}{0}
\setlength{\parindent}{0pt}
\setlength{\parskip}{5pt}
\setlength{\emergencystretch}{3em}
\newcommand{\sref}[2]{\hyperlink{source:#1:#2}{[S#2]}}
\newcommand{\cataloguescope}[1]{\paragraph{Recorded target.}#1}
\newcommand{\R}{\mathbb R}
\newcommand{\C}{\mathbb C}
\newcommand{\N}{\mathbb N}
\newcommand{\Z}{\mathbb Z}
\newcommand{\Q}{\mathbb Q}
\newcommand{\F}{\mathbb F}
\newcommand{\eps}{\varepsilon}
\begin{document}
\section{Linear Entropy and Rank Inequalities for Quantum Marginals}\label{statement}

\subsection{Definitions and mathematical statement}

\paragraph{Shared notation.}$\mathbb N=\{1,2,\ldots\}$; the standard number fields have their usual meanings. All additional conventions are stated locally.


\paragraph{Definitions and conventions.}
For a finite-dimensional density operator $\rho$, the marginal on subsystem $X$ is $\rho_X=\operatorname{Tr}_{X^c}\rho$. Its von Neumann entropy is $S(X)=-\operatorname{Tr}(\rho_X\log\rho_X)$, with $0\log0=0$, and its rank is $r_X=\operatorname{rank}\rho_X$. Linear entropy inequalities are linear combinations of the values $S(X)$; log-rank inequalities are linear combinations of $\log r_X$. The source\textquotesingle s entropy subquestion concerns Yeung--Zhang inequalities and their higher-order extensions, originally known for Shannon entropy.

\paragraph{Target retained from the canonical source catalogue.}
\paragraph{Catalogue wording retained for traceability.}
For a positive semidefinite quantum state $\rho$ on $A\otimes B\otimes C\otimes D$, determine the linear inequalities constraining the von Neumann entropies and logarithms of ranks of its marginals; in particular, decide the proposed higher-order entropy inequalities and whether $r_{AB}\leq r_{AC}r_{BC}$ for tripartite reductions.

\paragraph{Source scope and qualification.}
The named higher-order entropy family is source-defined. The rank inequality concerns a tripartite reduced state and is not an inequality between arbitrary unrelated matrices. Fresh status review of each component has not been established by this statement preparation.
\subsection{Short English statement}
Find the linear constraints on quantum marginal entropies and ranks, including the named entropy family and the indicated tripartite rank bound.
\subsection{Research attempt}
\paragraph{Model, date and investigated scope.}
OpenAI GPT-6 Astra (configured model \texttt{gpt-6-astra}), ultra reasoning; generated 11 October 2026. Proved the named rank inequality for pure tripartite states by support and coefficient matrices; does not resolve arbitrary reductions or the entropy families.

\paragraph{Sources and literature check.}
Read the supplied definition, catalogue wording, source qualifications and reference entries for this question. The historical source/context and inherited status are retained; no fresh literature or novelty certification is claimed. No external theorem is used without proof in this attempt; explicit model assumptions are distinguished from facts about the source model.

\paragraph{A pure-state rank test.} Let \(|\psi\rangle\) be a unit vector on finite-dimensional \(A\otimes B\otimes C\). If \(P_A,P_B\) project onto the supports of its marginals, then
\[
 \|((I-P_A)\otimes I)\psi\|^2=\operatorname{Tr}((I-P_A)\rho_A)=0,
\]
and similarly for \(B\). Hence \(\rho_{AB}\) is supported in \(\operatorname{ran}P_A\otimes\operatorname{ran}P_B\), giving \(r_{AB}\le r_A r_B\). Across a bipartition the coefficient matrix of a pure vector yields reduced matrices \(TT^*\) and the transpose of \(T^*T\); they have the same nonzero eigenvalues, by mapping a positive-eigenvalue vector with \(T\) or \(T^*\). Therefore
\[
 r_{AB}=r_C,\qquad r_{AC}=r_B,\qquad r_{BC}=r_A,
\]
so the source inequality \(r_{AB}\le r_{AC}r_{BC}\) holds in this pure tripartite subcase.
\paragraph{Boundary and scope.} All ranks are at least one for normalized states; product vectors give equality. The same coefficient-matrix argument gives equal complementary von Neumann entropies for pure states, but does not supply the named Yeung--Zhang family. A tripartite reduction of a four-party mixed or entangled state need not be pure, so complement equalities used above cannot simply be inserted into the full problem. The supplied statement records existing progress on the rank question; this elementary subcase is not claimed as new. A general mixed-state argument and the full four-party entropy/log-rank region remain outside this attempt.

\paragraph{Verification and adversarial audit.}
Checked support reduction, equality of nonzero coefficient-matrix spectra, rank-one boundaries and failure of the purity assumption for general reductions. This is a GPT-6 Astra Ultra self-audit, not independent human verification. Prompt A drafting followed by prompt B scope/dependency/boundary self-audit. Checked support reduction, equality of nonzero coefficient-matrix spectra, rank-one boundaries and failure of the purity assumption for general reductions. The remaining gap is: Handle mixed tripartite ranks without pure-state complement equalities and determine the full four-party entropy/log-rank constraints. Unresolved status and zero original-target completion are retained.

Coordinator checked marginal support and complementary pure-state spectra, confirming the rank inequality solely for the stated pure-tripartite subcase. This additional check was model-assisted, not independent human verification.

\paragraph{Tools, process and contributions.}
Python 3 and local file-reading/writing tools for manuscript serialization and structural checks; mathematical calculations were symbolic reasoning, with no numerical experiment or theorem prover. Contribution prompt A was followed by prompt B. GPT-6 Astra with ultra reasoning drafted the research and performed its stated self-audit; any additional coordinator contribution is identified in the verification and audit. A GPT-6-based Codex coordinator assembled the manuscript and checked repository integration; its exact variant and reasoning configuration were not exposed. No human mathematical contribution, independent human proof check or formal proof certification is supplied. Complete derivations and the consulted source references are included here; private reasoning transcripts are not published.

\paragraph{Final status and first remaining gap.}
\textbf{UNRESOLVED}. Handle mixed tripartite ranks without pure-state complement equalities and determine the full four-party entropy/log-rank constraints.

\section{Completion Estimate}
\noindent\textbf{COMPLETION: 0\%}
The proved calculation is an elementary auxiliary result or a test of a restricted strategy; no new portion of the full source target is established.

\subsection{Sources}
\begin{itemize}
\item[\hypertarget{source:statement:1}{S1}] ULAM.AI and the UnsolvedMath Contributors, UnsolvedMath, catalogue entry. Catalogue curation under CC BY 4.0; definitions and English exposition added locally.
\url{https://huggingface.co/datasets/ulamai/UnsolvedMath}.
\item[\hypertarget{source:statement:2}{S2}] Interactions between Operator Space Theory and Quantum Probability with Applications to Quantum Information (2018), Oberwolfach Reports, DOI 10.4171/owr/2018/22, pp. 1293--1379. Primary report accessed 2026-10-11; relevant PDF pages 82.
\url{https://doi.org/10.4171/owr/2018/22}.
\item[\hypertarget{source:statement:3}{S3}] Ranks of multipartite quantum states (catalogue-listed context; not independently consulted in this preparation).
\url{https://doi.org/10.1016/j.laa.2014.03.035}.
\item[\hypertarget{source:statement:4}{S4}] Four party zero entropy inequalities (catalogue-listed context; not independently consulted in this preparation).
\url{https://arxiv.org/abs/2105.07679}.
\end{itemize}
\label{end:statement}
\end{document}
